% !TEX root = ../main.tex
\chapter{Mockup e proposta di interazione}
\label{cap:mockup}
\bozza{Presentare i mockup come complemento visivo all'analisi dei requisiti,
collegando ogni schermata agli ID dei casi d'uso e ai passi dei template.
Selezionare poche schermate significative, ad esempio catalogo e dettaglio evento,
prenotazione e gestione amministrativa dell'evento.
Rappresentare anche un'alternativa rilevante, come la capienza esaurita.
Specificare che i mockup sono una proposta progettuale: non dimostrano
l'esistenza di una GUI, di sessioni applicative o di pagamenti integrati.}
\section{Navigazione proposta tra le schermate}
\bozza{Se utile, accompagnare i mockup con uno schema essenziale dei passaggi
fra le schermate; mantenere coerenti etichette, attori e riferimenti ai casi d'uso.}
